\subsection{Các kỹ thuật bổ trợ hiện có để giải quyết bài toán}

Trong bài toán phát hiện bất thường phụ tải điện, hiệu quả của mô hình không chỉ phụ thuộc vào thuật toán học máy hay học sâu cốt lõi mà còn chịu ảnh hưởng mang tính quyết định bởi cách thức biểu diễn dữ liệu và các kỹ thuật xử lý chuyên biệt. Hành vi tiêu thụ điện năng luôn gắn chặt với bối cảnh vật lý công trình, điều kiện khí tượng khu vực, tính chu kỳ sinh hoạt theo thời gian và hiện tượng mất cân bằng nhãn sâu sắc. Do đó, mục này tiến hành tổng quan các kỹ thuật bổ trợ hiện có trong tài liệu nghiên cứu, bao gồm kỹ thuật tương tác không gian vĩ mô, chuẩn hóa phụ tải theo ngữ cảnh, trích xuất động học chuỗi thời gian, phân cụm hành vi phụ tải, các hàm mất mát tùy biến và các cơ chế mô hình hóa nâng cao.

\subsubsection{Kỹ thuật tương tác giữa địa điểm và công năng}

Trong dữ liệu quản lý công trình, phụ tải tiêu thụ có sự phân hóa rõ rệt theo vị trí địa lý, điều kiện khí hậu và công năng sử dụng đặc thù của từng tòa nhà~\cite{bdg2, ashrae_gep3}. Trong bộ dữ liệu ASHRAE, thuộc tính \texttt{site\_id} định danh khu vực thu thập dữ liệu đi kèm trạm thời tiết tương ứng, trong khi \texttt{primary\_use} đại diện cho loại hình công năng hoạt động.

Để nắm bắt quy mô vận hành ở cấp độ vĩ mô, các kỹ thuật kỹ nghệ đặc trưng thường trích xuất diện tích sàn trung bình của các tòa nhà theo từng phân vùng địa lý và theo từng nhóm công năng:
\begin{equation*}
    \bar{A}_{S} = \mathbb{E}_{\mathbf{x} \in \mathcal{D}_S}[\texttt{square\_feet}], \qquad \bar{A}_{U} = \mathbb{E}_{\mathbf{x} \in \mathcal{D}_U}[\texttt{square\_feet}]
\end{equation*}
Trên cơ sở hai đại lượng thống kê này, đặc trưng tương tác không gian vĩ mô được xác định thông qua tích số giữa hai cấp độ diện tích trung bình:
\begin{equation*}
    \texttt{zone\_interaction} = \bar{A}_S \times \bar{A}_U
\end{equation*}
Nhằm tuân thủ nguyên tắc không rò rỉ dữ liệu, các giá trị kỳ vọng diện tích $\bar{A}_S$ và $\bar{A}_U$ bắt buộc phải được tính toán hoàn toàn trên tập dữ liệu huấn luyện, sau đó ánh xạ sang tập kiểm định và tập kiểm tra Out-of-sample.

Ưu điểm của kỹ thuật này là cung cấp cho các thuật toán phân lớp một thước đo định lượng về quy mô tương đối của công trình trong mối tương quan giữa khu vực và công năng sử dụng, hỗ trợ mô hình nhận diện các cụm phụ tải đặc thù mà không làm gia tăng số chiều phân loại. Đối với các công trình mới, thông tin siêu dữ liệu này giúp mô hình duy trì khả năng suy diễn khái quát hóa tốt hơn.

Tuy nhiên, nhược điểm là biến tương tác liên tục có độ phân tán giá trị rất lớn, đòi hỏi các thuật toán nhạy cảm với thang đo phải trải qua bước chuẩn hóa phù hợp. Ngoài ra, đây là các đặc trưng tĩnh ở cấp độ vĩ mô, do đó không thể nắm bắt được các biến động phụ tải vi mô tức thời theo từng khung giờ trong ngày của riêng từng tòa nhà.

\subsubsection{Chuẩn hóa phụ tải theo ngữ cảnh}

Lấy cảm hứng từ phương pháp chuẩn hóa phụ tải theo từng công trình của Orvati Nia và cộng sự~\cite{orvati2026}, phụ tải được biểu diễn tương đối so với mức tiêu thụ quen thuộc của chính công trình đó nhằm tạo ra bốn biến ngữ cảnh vi mô:
\begin{enumerate}
    \item Biến \texttt{z\_meter\_local} biểu diễn độ lệch chuẩn hóa của phụ tải so với mức tiêu thụ trung bình và độ lệch chuẩn của cùng công trình tại cùng khung giờ:
    \begin{equation*}
        z_{\text{local}}(t) = \frac{x_t - \mu_{b,h}}{\sigma_{b,h}}
    \end{equation*}
    trong đó $b$ đại diện cho mã công trình, $h$ là giờ trong ngày, còn $\mu_{b,h}$ và $\sigma_{b,h}$ được ước lượng từ tập huấn luyện.
    \item Biến \texttt{z\_meter\_workday} đo lường độ lệch phụ tải tương đối theo cùng công trình, cùng khung giờ và phân biệt rõ ngày làm việc với ngày nghỉ:
    \begin{equation*}
        z_{\text{workday}}(t) = \frac{x_t - \mu_{b,d,h}}{\sigma_{b,d,h}}
    \end{equation*}
    với $d$ là chỉ báo phân biệt ngày làm việc và ngày cuối tuần/ngày lễ.
    \item Biến \texttt{z\_meter\_robust} áp dụng kỹ thuật chuẩn hóa bền vững dựa trên khoảng tứ phân vị nhằm triệt tiêu ảnh hưởng của các giá trị ngoại lai cực trị:
    \begin{equation*}
        z_{\text{robust}}(t) = \frac{x_t - \operatorname{median}_{b}}{\operatorname{IQR}_{b}}
    \end{equation*}
    \item Biến \texttt{temp\_load\_inter} mô hình hóa sự tương tác trực tiếp giữa mức độ lệch phụ tải theo loại ngày và nhiệt độ môi trường xung quanh:
    \begin{equation*}
        \texttt{temp\_load\_inter}(t) = z_{\text{workday}}(t) \times \texttt{air\_temperature}(t)
    \end{equation*}
\end{enumerate}

Ưu điểm nổi bật của phương pháp chuẩn hóa theo ngữ cảnh là giúp mô hình phát hiện được các điểm tiêu thụ bất thường của một tòa nhà ngay cả khi giá trị công suất tuyệt đối vẫn nằm hoàn toàn trong dải biến thiên thông thường của toàn bộ quần thể dữ liệu. Sự kết hợp giữa thời gian biểu và nhiệt độ môi trường cung cấp bối cảnh vận hành chân thực để phân biệt giữa nhu cầu sử dụng đột biến do thời tiết với các sự cố bất thường thực sự.

Dù vậy, nhược điểm là kỹ thuật này phụ thuộc chặt chẽ vào độ tin cậy của chuỗi thống kê lịch sử. Nếu giai đoạn đo đạc cơ sở chứa quá nhiều dữ liệu lỗi hoặc độ dài quan trắc không đủ lớn, các giá trị kỳ vọng sẽ bị sai lệch. Hơn nữa, phương pháp này khó có thể áp dụng ngay cho các tòa nhà hoàn toàn mới khi chưa tích lũy đủ dữ liệu vận hành lịch sử.

\subsubsection{Đặc trưng động học chuỗi thời gian và thành phần phân rã}

Trong giải pháp vô địch cuộc thi LEAD, Fu và cộng sự~\cite{lead_winner} đã chứng minh vai trò quyết định của các đặc trưng động học chuỗi thời gian dưới dạng biến trễ, sai phân bậc một bậc hai và tỷ số biến thiên phụ tải. Dựa trên cơ sở đó, hệ thống tám đặc trưng động học mở rộng được thiết lập bao gồm:
\begin{enumerate}
    \item Các biến trễ vi mô \texttt{meter\_lag\_1h} và \texttt{meter\_lag\_2h} phản ánh mức tiêu thụ điện năng tại một và hai bước thời gian trước đó:
    \begin{equation*}
        \texttt{meter\_lag\_k}_t = m_{t-k}, \qquad k \in \{1,2\}
    \end{equation*}
    \item Các biến sai phân vận tốc \texttt{meter\_diff\_1h} và \texttt{meter\_diff\_2h} đo lường tốc độ thay đổi tức thời của phụ tải:
    \begin{equation*}
        \texttt{meter\_diff\_k}_t = m_t - m_{t-k}, \qquad k \in \{1,2\}
    \end{equation*}
    \item Các biến tỷ số động năng \texttt{meter\_ratio\_1h} và \texttt{meter\_ratio\_2h} nắm bắt biên độ gia tăng công suất tương đối:
    \begin{equation*}
        \texttt{meter\_ratio\_k}_t = \frac{m_t + 1}{m_{t-k} + 1}, \qquad k \in \{1,2\}
    \end{equation*}
    \item Biến thời lượng đóng băng \texttt{flatline\_duration} đếm số giờ liên tiếp mà mức tiêu thụ điện của công trình gần như không đổi ($|m_t - m_{t-1}| \leq 10^{-4}$), hỗ trợ nhận diện các sự cố cảm biến bị treo hoặc lỗi truyền thông.
    \item Biến thành phần phân rã \texttt{stl\_residual} đo lường độ sai lệch giữa phụ tải hiện tại so với trung vị của cùng khung giờ trong bảy ngày trước đó:
    \begin{equation*}
        \texttt{stl\_residual}_t = m_t - \operatorname{median}\{m_{t-24d} \mid d=1,\ldots,7\}
    \end{equation*}
\end{enumerate}

Ưu điểm lớn nhất của nhóm đặc trưng này là cung cấp tín hiệu trực tiếp về chiều hướng và tốc độ biến thiên phụ tải, cho phép các mô hình phát hiện tức thì các cú sốc tải hoặc sụt áp đột ngột. Cơ chế đếm đoạn phẳng giải quyết triệt để bài toán nhận diện cảm biến đóng băng, trong khi thành phần phần dư chu kỳ giúp loại bỏ nhịp điệu sinh hoạt 24 giờ để làm nổi bật dị thường.

Ngược lại, nhược điểm là các biến trễ đòi hỏi tính liên tục nghiêm ngặt của chuỗi thời gian. Tại các mốc khởi đầu chuỗi hoặc sau các khoảng trống dữ liệu bị mất kết nối, giá trị trễ sẽ bị khuyết và cần chiến lược điền bù dự phòng an toàn. Ngoài ra, ngưỡng phát hiện đoạn phẳng đòi hỏi phải được tinh chỉnh phù hợp với độ phân giải của từng loại công tơ đo đạc.

\subsubsection{Kỹ thuật phân cụm hành vi phụ tải}

Kỹ thuật phân cụm nhóm các tòa nhà dựa trên hình thái tiêu thụ thực tế thay vì phụ thuộc hoàn toàn vào nhãn công năng danh nghĩa. Mỗi tòa nhà có thể được biểu diễn bằng một vector đặc trưng hồ sơ phụ tải $\mathbf{v}_i \in \mathbb{R}^{D}$, chẳng hạn vector 168 chiều mô tả công suất tiêu thụ trung bình qua từng giờ trong tuần. Thuật toán K-Means~\cite{kmeans} tối ưu hóa hàm mục tiêu phân cụm theo khoảng cách:
\begin{equation*}
    \mathcal{J} = \sum_{k=1}^{K} \sum_{\mathbf{v}_i \in \mathcal{C}_k} \left\| \mathbf{v}_i - \boldsymbol{\mu}_k \right\|_2^2
\end{equation*}
với $\boldsymbol{\mu}_k$ là vector trọng tâm của cụm thứ $k$. Phương pháp này từng được Cansever và Konyar~\cite{cansever2026} ứng dụng để chia nhóm công trình trước khi huấn luyện các mô hình phát hiện dị thường chuyên biệt.

Về ưu điểm, phân cụm hồ sơ tiêu thụ giúp phân tách các nhóm công trình có hình thái phụ tải tương đồng, phản ánh trung thực thực tế vận hành vượt ra ngoài ranh giới phân loại công năng ban đầu. Nhãn cụm có thể tích hợp trực tiếp vào mô hình như một đặc trưng bổ trợ giá trị.

Tuy vậy, nhược điểm là kết quả phân cụm rất nhạy cảm với việc lựa chọn số lượng cụm $K$ và độ đo khoảng cách. Khoảng cách Euclidean thông thường dễ bị ảnh hưởng bởi hiện tượng lệch pha thời gian hoặc chênh lệch quy mô công suất giữa các tòa nhà lớn và nhỏ.

\subsubsection{Hàm mất mát entropy chéo có trọng số}

Nhằm xử lý trực tiếp tình trạng mất cân bằng lớp nghiêm trọng trong bài toán phát hiện dị thường, hàm mất mát entropy chéo có trọng số gán hệ số phạt $w_+$ cao hơn cho các sai sót phân loại trên lớp bất thường:
\begin{equation*}
    \mathcal{L}_{\text{WBCE}} = -\frac{1}{N}\sum_{i=1}^{N}\left[ w_{+}\, y_i \log \hat{p}_i + (1 - y_i)\log\left(1 - \hat{p}_i\right)\right]
\end{equation*}
với $\hat{p}_i$ là xác suất dự đoán và $y_i \in \{0, 1\}$. Hệ số này tương ứng với tham số điều chỉnh tỷ lệ mẫu dương trong các thư viện học máy và học sâu phổ biến.

Ưu điểm của kỹ thuật là tính đơn giản, dễ triển khai và có thể tích hợp trực tiếp vào quá trình tối ưu hóa đạo hàm của mô hình mà không cần can thiệp lấy mẫu lại dữ liệu, từ đó ngăn chặn hiện tượng mô hình hội tụ về giải pháp tầm thường là dự đoán toàn bộ mẫu thuộc lớp bình thường.

Mặc dù vậy, việc áp đặt trọng số phạt quá cao có thể làm méo mó xác suất dự đoán thực tế, khiến mô hình phản ứng quá nhạy với các biến động phụ tải nhẹ và làm gia tăng đáng kể số lượng cảnh báo giả, đòi hỏi phải tiến hành dò tìm lại ngưỡng phân loại tối ưu.

\subsubsection{Hàm mất mát Focal Loss}

Hàm mất mát Focal Loss được Lin và cộng sự~\cite{focal_loss} phát triển nhằm giải quyết triệt để sự áp đảo của các mẫu dễ học trong bài toán mất cân bằng nhãn. Bằng cách bổ sung thừa số suy giảm động $(1 - p_t)^\gamma$, hàm mất mát tự động triệt tiêu đóng góp đạo hàm từ các mẫu mà mô hình đã phân loại tự tin:
\begin{equation*}
    \mathrm{FL}(p_t) = -\alpha_t \left(1 - p_t\right)^{\gamma} \log\left(p_t\right)
\end{equation*}
trong đó $p_t$ là xác suất mô hình gán cho lớp chính xác, $\alpha_t$ là hệ số cân bằng nhãn và $\gamma \geq 0$ là tham số điều biến mức độ tập trung vào các mẫu khó.

Ưu điểm nổi bật của Focal Loss là tập trung tối đa tài nguyên học của mô hình vào các trường hợp biên khó phân định, hạn chế tình trạng hàm mất mát bị chi phối bởi hàng triệu quan trắc bình thường dễ nhận biết.

Nhược điểm là hàm đưa thêm hai siêu tham số $(\alpha, \gamma)$ vào không gian tìm kiếm, đòi hỏi chi phí tinh chỉnh bổ sung. Ngoài ra, điểm số đầu ra của mô hình sau khi tối ưu với Focal Loss không còn mang ý nghĩa xác suất hiệu chuẩn thực tế và cần qua bước cân chỉnh nếu muốn sử dụng làm xác suất tin cậy.

\subsubsection{Mạng tích chập thời gian kết hợp bộ tự mã hóa TCN}

Kiến trúc TCN kết hợp bộ tự mã hóa ứng dụng các phép tích chập nhân quả giãn cách trong cả nhánh mã hóa và giải mã nhằm học cách khôi phục chuỗi thời gian thông qua một nút thắt cổ chai đại diện~\cite{tcn, handigol2025}. Điểm số bất thường được xác định dựa trên mức độ sai số tái tạo giữa chuỗi thực tế và chuỗi khôi phục.

Ưu điểm của mô hình là tốc độ huấn luyện vượt bậc nhờ khả năng song song hóa hoàn toàn của phép toán tích chập một chiều trên GPU, kết hợp cùng vùng tiếp nhận mở rộng lũy thừa giúp bao quát toàn diện các chu kỳ phụ tải ngày và tuần.

Tuy nhiên, nhược điểm là nếu không gian nút thắt cổ chai không được thiết kế đủ hẹp, mạng có thể học tái tạo cả các đoạn phụ tải bất thường. Bên cạnh đó, các phép toán tích chập nhiều kênh trên kích thước lô lớn đòi hỏi dung lượng bộ nhớ GPU đáng kể.

\subsubsection{Cơ chế tự chú ý trong mô hình chuỗi}

Cơ chế tự chú ý trong kiến trúc Transformer~\cite{transformer} cho phép tính toán trực tiếp trọng số tương quan giữa mọi cặp thời điểm trong chuỗi phụ tải thông qua các phép chiếu ma trận truy vấn, khóa và giá trị:
\begin{equation*}
    \text{Attention}(\mathbf{Q}, \mathbf{K}, \mathbf{V}) = \text{softmax}\left( \frac{\mathbf{Q}\mathbf{K}^\top}{\sqrt{d_k}} \right) \mathbf{V}
\end{equation*}

Ưu điểm chính của cơ chế tự chú ý là khả năng mô hình hóa sự phụ thuộc xa mà không bị suy giảm thông tin theo bước truyền như mạng hồi quy, cho phép liên kết trực tiếp giữa các thời điểm cách nhau nhiều ngày trong cùng một chu kỳ vận hành. Ma trận chú ý đồng thời cung cấp khả năng trực quan hóa mức độ quan trọng của các mốc thời gian lịch sử đối với dự đoán hiện tại.

Dù vậy, nhược điểm là chi phí tính toán và bộ nhớ tăng theo hàm bậc hai $O(w^2)$ theo chiều dài cửa sổ trượt, gây hạn chế lớn khi muốn mở rộng phạm vi ngữ cảnh. Hơn nữa, cơ chế này đòi hỏi lượng dữ liệu chuỗi rất lớn để tối ưu hóa ổn định, dễ rơi vào trạng thái quá khớp khi tập dữ liệu bị giới hạn về số lượng công trình.

\subsubsection{Hàm mất mát căn chỉnh thời gian mềm khả vi Soft-DTW}

Kỹ thuật Soft-DTW~\cite{soft_dtw} là phiên bản làm trơn khả vi của thuật toán căn chỉnh thời gian động cổ điển, từng được Prabhu và cộng sự~\cite{prabhu2024} sử dụng làm hàm mất mát tái tạo cho mô hình phát hiện dị thường chuỗi năng lượng:
\begin{equation*}
    \text{Soft-DTW}_\gamma(\mathbf{x}, \hat{\mathbf{x}}) = \operatorname{soft-min}_\gamma \left( \left\{ \langle \mathbf{A}_\pi, \boldsymbol{\Delta}(\mathbf{x}, \hat{\mathbf{x}}) \rangle \right\}_{\pi \in \mathcal{A}} \right)
\end{equation*}
với $\boldsymbol{\Delta}$ là ma trận khoảng cách giữa các điểm chuỗi, $\mathcal{A}$ là tập hợp các đường căn chỉnh hợp lệ và phép lấy cực tiểu mềm $\operatorname{soft-min}_\gamma$ cho phép tính toán đạo hàm phục vụ lan truyền ngược.

Ưu điểm của Soft-DTW là khả năng dung nạp các sai lệch lệch pha thời gian nhỏ trong nhịp phụ tải, giúp mô hình không phạt nặng các biến động tiêu thụ bình thường chỉ bị trễ giờ nhẹ so với lịch biểu thông thường.

Ngược lại, nhược điểm là thuật toán quy hoạch động có độ phức tạp tính toán bậc hai $O(w^2)$, làm chậm đáng kể quá trình huấn luyện mô hình. Đặc biệt, nếu sự cố bất thường trong thực tế chính là việc chậm giờ đóng cắt thiết bị, cơ chế co dãn đàn hồi của Soft-DTW có thể vô tình xem sự cố này là hành vi bình thường và bỏ sót lỗi.

\subsubsection{Mô hình chuỗi có điều kiện theo siêu dữ liệu}

Hướng tiếp cận có điều kiện kết hợp các thuộc tính siêu dữ liệu tĩnh như vị trí khu vực, loại hình công năng hoặc nhãn cụm hành vi vào không gian biểu diễn ẩn của mạng chuỗi thông qua cơ chế điều biến tuyến tính FiLM~\cite{film}:
\begin{equation*}
    \mathbf{h}_t' = \boldsymbol{\gamma}(\mathbf{e}) \odot \mathbf{h}_t + \boldsymbol{\beta}(\mathbf{e})
\end{equation*}
trong đó $\mathbf{e}$ là vector nhúng của siêu dữ liệu, $\mathbf{h}_t$ là trạng thái ẩn của chuỗi thời gian, còn $\boldsymbol{\gamma}$ và $\boldsymbol{\beta}$ là các mạng biến đổi tuyến tính điều khiển hệ số co giãn và độ lệch.

Ưu điểm của kỹ thuật là giúp mạng thích ứng linh hoạt các trạng thái biểu diễn chuỗi theo bối cảnh vật lý cụ thể của từng công trình, nâng cao khả năng khái quát hóa khi áp dụng mô hình chung cho nhiều nhóm tòa nhà khác nhau.

Tuy nhiên, việc tích hợp thêm các ma trận nhúng và mạng điều biến làm gia tăng số lượng tham số, đòi hỏi các biện pháp điều chuẩn chặt chẽ để phòng chống quá khớp. Hơn nữa, mô hình phụ thuộc nghiêm ngặt vào tính sẵn sàng của siêu dữ liệu tại thời điểm suy luận.

\subsubsection{Mô hình sinh đối kháng kết hợp mạng LSTM}

Mô hình GAN-LSTM kết hợp mạng sinh đối kháng với các tầng hồi quy LSTM nhằm học phân phối của chuỗi phụ tải bình thường trên các cửa sổ thời gian kéo dài (Orvati Nia và cộng sự~\cite{orvati2026}). Khi suy luận, điểm bất thường kết hợp giữa sai số tái tạo và độ lệch đặc trưng phân biệt:
\begin{equation*}
    \mathcal{S}(\mathbf{x}) = \lambda\,\mathcal{L}_{\mathrm{rec}}\!\left(\mathbf{x}, G(\mathbf{z}^*)\right) + (1-\lambda)\,\mathcal{L}_{\mathrm{feat}}\!\left(f_D(\mathbf{x}), f_D(G(\mathbf{z}^*))\right)
\end{equation*}
trong đó $\mathbf{z}^*$ là vector tiềm ẩn tối ưu được tìm kiếm tại thời điểm suy luận và $f_D$ là biểu diễn đặc trưng trung gian của mạng phân biệt.

Ưu điểm của phương pháp là học được phân phối chuẩn xác của chuỗi thời gian mà không phụ thuộc vào nhãn bất thường, đồng thời tận dụng bộ nhớ LSTM để duy trì hình thái tuần hoàn của phụ tải qua nhiều khung giờ.

Mặt khác, nhược điểm lớn là quá trình tối ưu hóa vector tiềm ẩn khi suy luận đòi hỏi hàng trăm vòng lặp tính toán cho mỗi cửa sổ, tạo ra nút thắt cổ chai nghiêm trọng về tốc độ xử lý khi áp dụng cho các hệ thống giám sát quy mô lớn. Việc đánh giá theo cửa sổ gộp cũng làm giảm độ chính xác định vị thời điểm phát sinh bất thường.

\subsubsection{Mô hình sinh đối kháng tích chập một chiều}

Kiến trúc DCGAN một chiều sử dụng các tầng tích chập sâu để xử lý cửa sổ phụ tải điện năng (Handigol và Arjunan~\cite{handigol2025}). Mạng tạo sinh sinh ra chuỗi nhân tạo từ vector ngẫu nhiên, trong khi mạng phân biệt đối kháng đánh giá tính chân thực của chuỗi đo đạc. Điểm bất thường được suy luận trực tiếp từ sai số tái tạo:
\begin{equation*}
    \mathcal{S}(\mathbf{x}) = \mathcal{L}_{\mathrm{rec}}(\mathbf{x}, \hat{\mathbf{x}})
\end{equation*}

Ưu điểm nổi bật là tốc độ tính toán song song nhanh hơn đáng kể so với các mô hình GAN dựa trên LSTM, cùng khả năng trích xuất nhạy bén các đặc trưng hình thái sóng cục bộ như gai tải đột biến hoặc các khoảng giá trị đóng băng.

Dù vậy, việc sử dụng các tầng tích chập chuyển vị có thể gây ra hiện tượng méo chuỗi hoặc sinh mẫu tuần hoàn giả nếu không thiết lập cẩn trọng kích thước bộ lọc. Ngoài ra, khả năng duy trì sự phụ thuộc dài hạn của mạng tích chập bị giới hạn bởi độ sâu của các tầng tích chập.

\subsubsection{Kỹ thuật tái tạo song song nhiều điểm}

Để giải quyết nút thắt cổ chai về tốc độ của mô hình AnoGAN~\cite{anogan}, kỹ thuật tái tạo song song thực hiện tối ưu hóa đồng thời nhiều cửa sổ quan sát theo lô trên bộ xử lý đồ họa GPU (Prabhu và cộng sự~\cite{prabhu2024}):
\begin{equation*}
    \mathbf{z}^* = \arg\min_{\mathbf{z}} \mathcal{L}_{\mathrm{rec}}\!\left(\mathbf{x}, G(\mathbf{z})\right)
\end{equation*}

Ưu điểm của kỹ thuật là tận dụng năng lực xử lý đa luồng của phần cứng để tăng thông lượng suy luận, giảm thiểu thời gian chờ đợi khi phải rà soát các tệp dữ liệu lớn.

Tuy nhiên, bản chất của phương pháp vẫn đòi hỏi các bước cập nhật độ dốc lặp đi lặp lại tại thời điểm suy luận, do đó mức độ tiêu tốn bộ nhớ GPU vẫn còn tương đối cao và tốc độ chưa thể so sánh trực tiếp với các mô hình phân lớp xuôi chiều một lượt.

Tổng thể các kỹ thuật bổ trợ trên cung cấp một hệ thống giải pháp đa tầng từ cấp độ không gian vĩ mô, ngữ cảnh vi mô cho đến động học chuỗi thời gian, tạo nền tảng lý thuyết vững chắc để nhóm nghiên cứu xây dựng và tối ưu hóa các không gian dữ liệu thực nghiệm trong các chương tiếp theo.
